% !TEX root = ./adlt_homework_8.tex
\documentclass[12pt]{article}

% --- Core math + proof tools ---
\usepackage{amsmath, amssymb, amsthm}

% --- Lists and spacing ---
\usepackage{enumitem}
\setlist[enumerate,1]{label=\textbf{(\alph*)}, leftmargin=2em, itemsep=0.25em}
\setlist[itemize]{leftmargin=1.5em, itemsep=0.25em}
\setlength{\parskip}{0.35em}
\setlength{\parindent}{0pt}

% --- Page geometry + header/footer ---
\usepackage{geometry}
\geometry{margin=1in}
\usepackage{csquotes} % to use blockquotes
\usepackage{graphicx} % for figures
\graphicspath{{images/}} % tells LaTeX to look in the “images” folder

\usepackage{fancyhdr}
\pagestyle{fancy}
\fancyhf{}
\rhead{MATH 421 — 003, Fall 2025}
\lhead{Homework 10}
\cfoot{\thepage}
\renewcommand{\headrulewidth}{0.4pt}

% --- Links (nice for references, stays subtle) ---

\usepackage[hidelinks]{hyperref}
\setcounter{tocdepth}{1}

% --- Colored definition/theorem boxes ---
\usepackage[most]{tcolorbox}

% Reusable styles for boxes
\tcbset{
  definitionstyle/.style={
    enhanced,
    colback=blue!3!white,
    colframe=blue!60!black,
    fonttitle=\bfseries,
    attach boxed title to top left={yshift=-2mm,xshift=2mm},
    boxed title style={colback=blue!60!black, colframe=blue!60!black},
    coltitle=white,
    rounded corners,
    shadow={0.7mm}{-0.7mm}{0mm}{black!20!white},
    breakable,
  },
  theoremstyle/.style={
    enhanced,
    colback=green!3!white,
    colframe=green!50!black,
    fonttitle=\bfseries,
    attach boxed title to top left={yshift=-2mm,xshift=2mm},
    boxed title style={colback=green!50!black, colframe=green!50!black},
    coltitle=white,
    rounded corners,
    shadow={0.7mm}{-0.7mm}{0mm}{black!20!white},
    breakable,
  }
}

% Environments you can use in discussions
\newenvironment{definitionbox}[1][]% optional title
{\begin{tcolorbox}[definitionstyle,title={#1}]}%
{\end{tcolorbox}}

\newenvironment{theorembox}[1][]% optional title
{\begin{tcolorbox}[theoremstyle,title={#1}]}%
{\end{tcolorbox}}

% --- Title info ---
\title{Homework 10}
\author{Alejandro De La Torre \\ MATH 421 — The Theory of Single Variable Calculus \\ Section 003 \\ Instructor: Grace Work}
\date{December 2025}

% --- proof environments that might be relevant ---
\newtheorem*{claim}{Claim}
\newtheorem*{lemma}{Lemma}
\newtheorem*{remark}{Remark}

% --- Convenience: a Problem header command (section-like with TOC/bookmarks) ---
% Usage: \problem[Optional short title]{<number>}
\makeatletter
\newcommand{\problem}[2][]{%
  \def\prob@title{Problem #2\if\relax\detokenize{#1}\relax\else\ (\textit{#1})\fi}%
  \phantomsection%
  \pdfbookmark[1]{\prob@title}{prob#2}%
  \addcontentsline{toc}{section}{\prob@title}%
  \section*{\prob@title}%
  \vspace{-0.25em}\hrule\vspace{0.8em}%
}
\makeatother

% --- Convenience: an ignore command to comment out blocks ---
\newcommand{\ignore}[1]{}


\begin{document}
\maketitle

% ------------------ collaboration + sources ------------------
\section*{Collaboration Statement}
% (List any classmates you collaborated with here, or write ``I worked alone'' if you did not collaborate with anyone.)
I worked on this homework independently. No outside collaborators or resources were used.

\tableofcontents
\bigskip
\hrule
\bigskip

\newpage


% ------------------ problems ------------------
\problem[Ross Chapter 3 18.5 (a)]{1}
Let $f$ and $g$ be continuous functions in $[a,b]$ such that $f(a) \geq g(a)$ and $f(b) \leq g(b)$. Prove $f(x_0) = g(x_0)$ for at least one $x_0$ in $[a,b]$.
\vspace{0.75em}

\textbf{Discussion.}\\
We want a point where $f(x_0) = g(x_0)$.  This is equivalent to $f(x_0) - g(x_0) = 0$, so it is natural to consider the difference $h(x) = f(x) - g(x)$.  Since $f$ and $g$ are continuous, $h$ is continuous on $[a,b]$.  The inequalities $f(a)\ge g(a)$ and $f(b)\le g(b)$ then say $h(a)\ge 0$ and $h(b)\le 0$, so $0$ lies between $h(a)$ and $h(b)$.  By the Intermediate Value Theorem applied to $h$, there is some $x_0\in [a,b]$ with $h(x_0)=0$, i.e.\ $f(x_0)=g(x_0)$.
\vspace{0.25em}

\textbf{Proof.}\\
Define $h : [a,b] \to \mathbb{R}$ by
\[
  h(x) = f(x) - g(x).
\]
Since $f$ and $g$ are continuous on $[a,b]$ and the difference of two continuous
functions is continuous, it follows that $h$ is continuous on $[a,b]$.

From the hypotheses we have
\[
  f(a) \ge g(a) \quad\Longrightarrow\quad h(a) = f(a) - g(a) \ge 0,
\]
and
\[
  f(b) \le g(b) \quad\Longrightarrow\quad h(b) = f(b) - g(b) \le 0.
\]
Thus $0$ lies between $h(a)$ and $h(b)$; equivalently, $h(b) \le 0 \le h(a)$.

By the \hyperref[thm:18.2]{Intermediate Value Theorem}, since $h$ is continuous
on $[a,b]$ and $0$ lies between $h(a)$ and $h(b)$, there exists some $x_0 \in [a,b]$
such that
\[
  h(x_0) = 0.
\]
By the definition of $h$ this means
\[
  f(x_0) - g(x_0) = 0,
\]
so $f(x_0) = g(x_0)$.

Therefore there exists at least one $x_0 \in [a,b]$ with $f(x_0) = g(x_0)$, as required. \qed

\newpage

\problem[Ross Chapter 3 18.10]{2}
Suppose $f$ is continuous on $[0,2]$ and $f(0)=f(2)$. Prove there exist $x,y$ in $[0,2]$ such that $|y-x|=1$ and $f(x)=f(y)$. \textit{Hint}: Consider $g(x) = f(x+1)-f(x)$ on $[0,1]$.
\vspace{0.75em}

\textbf{Discussion.} \\
We want two points $x,y$ with $|y-x|=1$ and $f(x)=f(y)$.  Writing $y=x+1$, this is
equivalent to finding $x\in[0,1]$ such that
\[
f(x+1)-f(x)=0.
\]
So it is natural to define
\[
g(x)=f(x+1)-f(x) \text{, for }x\in[0,1].
\]
Because $f$ is continuous on $[0,2]$, both $f(x+1)$ and $f(x)$ are continuous on $[0,1]$,
so their difference $g$ is continuous on $[0,1]$.  The condition $f(0)=f(2)$ then forces
\[
g(0)=f(1)-f(0),\qquad g(1)=f(2)-f(1)=f(0)-f(1)=-g(0),
\]
so either $g(0)=0$ or $g(0)$ and $g(1)$ have opposite signs.  In either case, $0$ lies between $g(0)$ and $g(1)$, so by the \hyperref[thm:18.2]{Intermediate Value Theorem} there is some $x\in[0,1]$ with $g(x)=0$, i.e.\ $f(x+1)=f(x)$ and $| (x+1)-x|=1$.

\textit{\textbf{TL;DR}}: We basically have to translate the statement we are trying to prove to make it in terms of just $x$ (which we do by defining $g$) to then make the question about applying the IVT to $g$ to prove that it has a zero, which then maps back to the original statement we are trying to prove about $f$.

\vspace{0.25em}

\textbf{Proof.}\\
Let $g$ be a function defined on $[0,1]$ as
\[
g(x)=f(x+1)-f(x).
\]
Since $f$ is continuous on $[0,2]$ (and $[0,1] \subset [0,2]$), then $g$ must be continuous on $[0,1]$.

We compute the endpoint values of $g$:
\[
g(0)=f(1)-f(0),
\]
and, using $f(0)=f(2)$,
\[
g(1)=f(2)-f(1)=f(0)-f(1)=-g(0).
\]

If $g(0)=0$, then taking $x=0$ and $y=1$ gives
\[
|y-x|=|1-0|=1,\qquad f(x)=f(0)=f(1)=f(y),
\]
so the desired property holds.

Otherwise $g(0)\neq 0$, so $g(1)=-g(0)$ has the opposite sign. Thus
\[
g(0)\cdot g(1)<0,
\]
which means $0$ lies strictly between $g(0)$ and $g(1)$. Since $g$ is continuous on $[0,1]$, the \hyperref[thm:18.2]{Intermediate Value Theorem} guarantees the existence of some $x\in(0,1)$ such that
\[
g(x)=0.
\]
For this $x$, set $y=x+1$.  Then $x,y\in[0,2]$, we have
\[
|y-x|=|x+1-x|=1,
\]
and
\[
f(y)=f(x+1)=f(x),
\]
because $g(x)=f(x+1)-f(x)=0$.

Therefore there exist $x,y\in[0,2]$ with $|y-x|=1$ and $f(x)=f(y)$, as required. \qed

\newpage

\problem[Ross Chapter 3 19.6 (b)]{3}
\textit{(From Part (a):) Let $f(x) = \sqrt{x}$ for $x \geq 0$. Show $f'$ is unbounded on $(0,1]$ but $f$ is nevertheless uniformly continuous on $(0,1]$ Compare with Theorem 19.6.}

(b) Show that $f$ is uniformly continuous on $[1,\infty)$.
\vspace{0.75em}

\textbf{Discussion.} \\
We want to prove that $f(x) = \sqrt{x}$ is \emph{uniformly} continuous on $[1,\infty)$.
By \hyperref[def:19.1]{Definition 19.1}, this means:

\[
\text{for each } \varepsilon > 0 \text{ there exists } \delta > 0 \text{ such that }
x,y \in [1,\infty),\ |x-y| < \delta \Longrightarrow |\sqrt{x} - \sqrt{y}| < \varepsilon.
\]

So the game is: fix an arbitrary $\varepsilon>0$ and then cook up a $\delta$ that depends
\emph{only} on $\varepsilon$, not on a particular choice of $x$ or $y$.

A standard trick for square roots is to use the algebraic identity
\[
|\sqrt{x}-\sqrt{y}|
 = \left|\frac{x-y}{\sqrt{x}+\sqrt{y}}\right|.
\]
On our domain $[1,\infty)$ we always have $\sqrt{x} \ge 1$ and $\sqrt{y} \ge 1$, so
\[
\sqrt{x} + \sqrt{y} \ge 2.
\]
That means we can bound the difference by
\[
|\sqrt{x}-\sqrt{y}| \le \frac{|x-y|}{2}.
\]

This is the key estimate: on $[1,\infty)$, the function $\sqrt{x}$ cannot change
faster than ``slope $\tfrac12$'' in the sense that
\[
|\sqrt{x}-\sqrt{y}| \le \frac12\,|x-y|.
\]
Once we have this, the $\varepsilon$–$\delta$ choice is easy: if we make
$|x-y|$ smaller than $2\varepsilon$, then $\frac12|x-y|$ will automatically be
smaller than $\varepsilon$. So in the proof we will take
\[
\delta = 2\varepsilon
\]
and then check that this choice satisfies the definition of uniform continuity
on $[1,\infty)$.
\vspace{0.25em}

\textbf{Proof.}\\
Let $\varepsilon > 0$, and $\delta = 2\varepsilon$. From \hyperref[def:19.1]{Definition 19.1} let $x,y \in S$ with $|x-y| < \delta$.
Then,
\[
|\sqrt{x} - \sqrt{y}|
  = \left|\frac{x-y}{\sqrt{x} + \sqrt{y}}\right|.
\]

Because $x,y \in [1,\infty)$, we have $\sqrt{x} \ge 1$ and $\sqrt{y} \ge 1$, so
\[
\sqrt{x} + \sqrt{y} \ge 2.
\]
Therefore
\[
|\sqrt{x} - \sqrt{y}|
  = \frac{|x-y|}{\sqrt{x} + \sqrt{y}}
  \le \frac{|x-y|}{2}.
\]

Using $|x-y| < \delta = 2\varepsilon$, we obtain
\[
|\sqrt{x} - \sqrt{y}|
  \le \frac{|x-y|}{2}
  < \frac{2\varepsilon}{2}
  = \varepsilon.
\]

Since $\varepsilon>0$ was arbitrary, by \hyperref[def:19.1]{Definition 19.1}, $f(x)=\sqrt{x}$ is uniformly continuous on $[1,\infty)$. 
\qed

\newpage

\problem[Ross Chapter 3 19.7]{4}
(a) Let $f$ be a continuous function on $[0, \infty)$. Prove that if $f$ is uniformly continuous on $[k, infty)$ for some $k$, then $f$ is uniformly continuous on $[0,\infty).$

(b) Use (a) and Exercise 19.6(b) to show that $f(x) = \sqrt{x}$ is uniformly continuous on $[0, \infty)$.
\vspace{0.75em}

\textbf{Discussion.} \\
%(Very brief plan / key ideas.)
\vspace{0.25em}

\textbf{Proof.}\\
% (Your proof goes here.)
\qed

\newpage

\problem[Uniform Continuity of $f$]{5}
%\textit{(Problem statement goes here if needed.)}
Suppose $f : \mathbb{R} \rightarrow \mathbb{R}$. Prove that if there exists a real number $M \geq 0$ such that $f(x) - f(y)| \leq M|x-y|$ for all $x,y \in \mathbb{R}$, then $f$ is uniformly continuous on $\mathbb{R}$.
\vspace{0.75em}

\textbf{Discussion.} \\
%(Very brief plan / key ideas.)
\vspace{0.25em}

\textbf{Proof.}\\
% (Your proof goes here.)
\qed

\newpage


\problem[Differntiability of $f$]{6}
Suppose that $0 \in (a,b)$ and $f(x)=xg(x)$ for some function $g: (a,b) \rightarrow \mathbb{R}$ which is continuous at $x=0$. Prove that $f$ is differentiable at $x=0$, and find $f'(0)$ in terms of $g$.

\textit{Hint: Note that you cannot use the product rule, since $g(x)$ is not guaranteed to be differentiable at $x=0$}

\vspace{0.75em}

\textbf{Discussion.} \\
%(Very brief plan / key ideas.)
\vspace{0.25em}

\textbf{Proof.}\\
% (Your proof goes here.)
\qed


\newpage
\appendix
\section*{Appendix: Theorems}
\addcontentsline{toc}{section}{Appendix: Theorems Used}

% ===============================
% 17.2 Theorem — ε–δ definition of continuity
% ===============================
\begin{theorembox}[Theorem 17.2]
\label{thm:17.2}
Let $f$ be a real-valued function whose domain is a subset of $\mathbb{R}$.  
Then $f$ is continuous at $x_{0} \in \mathrm{dom}(f)$ if and only if  
\[
\begin{gathered}
\text{for each } \varepsilon > 0 \text{ there exists } \delta > 0 \text{ such that } 
x \in \mathrm{dom}(f) \text{ and}\\[0.4em]
|x - x_{0}| < \delta \;\Longrightarrow\; |f(x) - f(x_{0})| < \varepsilon.
\end{gathered}
\]
\end{theorembox}

% ===============================
% 18.1 Theorem Extreme Value Theorem
% ===============================
\begin{theorembox}[Theorem 18.1 — Extreme Value Theorem]
\label{thm:18.1}
Let $f$ be a continuous real-valued function on a closed interval $[a,b]$.  
Then $f$ is a bounded function. Moreover, $f$ assumes its maximum and minimum
values on $[a,b]$; that is, there exist $x_{0}, y_{0}$ in $[a,b]$ such that  
\[
  f(x_{0}) \le f(x) \le f(y_{0}) \quad \text{for all } x \in [a,b].
\]
\end{theorembox}

% ===============================
% 18.2 Intermediate Value Theorem
% ===============================
\begin{theorembox}[Theorem 18.2 — Intermediate Value Theorem]
\label{thm:18.2}
If $f$ is a continuous real-valued function on an interval $I$, then $f$ has  
the intermediate value property on $I$: Whenever $a,b \in I$, $a < b$, and $y$  
lies between $f(a)$ and $f(b)$ (i.e., $f(a) < y < f(b)$ or $f(b) < y < f(a)$),  
there exists at least one $x$ in $(a,b)$ such that  
\[
  f(x) = y.
\]
\end{theorembox}

% ===============================
% 18.3 Corollary
% ===============================
\begin{theorembox}[Corollary 18.3]
\label{thm:18.3}
If $f$ is a continuous real-valued function on an interval $I$, then the set  
\[
  f(I) = \{ f(x) : x \in I \}
\]
is also an interval or a single point.
\end{theorembox}

% ===============================
% 19.1 Definition — Uniform Continuity
% ===============================
\begin{definitionbox}[Definition 19.1 — Uniform Continuity]
\label{def:19.1}
Let $f$ be a real-valued function defined on a set $S \subseteq \mathbb{R}$.  
Then $f$ is \emph{uniformly continuous} on $S$ if  

\[
\text{for each } \varepsilon > 0 \text{ there exists } \delta > 0 \text{ such that }  
x,y \in S \text{ and } |x-y| < \delta \;\Rightarrow\; |f(x)-f(y)| < \varepsilon.
\]

We will say $f$ is \emph{uniformly continuous} if $f$ is uniformly continuous  
on $\mathrm{dom}(f)$.
\end{definitionbox}

% ===============================
% 19.2 Theorem — Uniform Continuity on [a,b]
% ===============================
\begin{theorembox}[Theorem 19.2]
\label{thm:19.2}
If $f$ is continuous on a closed interval $[a,b]$,  
then $f$ is uniformly continuous on $[a,b]$.
\end{theorembox}


% ===============================
% 19.3 Discussion
% ===============================
\begin{definitionbox}[Discussion 19.3]
Example 5 illustrates the power of Theorem~\ref{thm:19.2}, but it still may  
not be clear why uniform continuity is worth studying. One of the important  
applications of uniform continuity concerns the integrability of continuous  
functions on closed intervals.

To see the relevance of uniform continuity, consider a continuous nonnegative  
real-valued function $f$ on $[0,1]$. For $n \in \mathbb{N}$ and  
$i = 0,1,2,\ldots,n-1$, let  
\[
M_{i,n} = \sup\{f(x) : x \in [\tfrac{i}{n}, \tfrac{i+1}{n}]\},  
\qquad
m_{i,n} = \inf\{f(x) : x \in [\tfrac{i}{n}, \tfrac{i+1}{n}]\}.
\]

Then the upper sum and lower sum are given by  
\[
U_n = \frac{1}{n} \sum_{i=0}^{n-1} M_{i,n},
\qquad
L_n = \frac{1}{n} \sum_{i=0}^{n-1} m_{i,n}.
\]

The function $f$ would be Riemann integrable provided $U_n - L_n \to 0$, i.e.,
\[
\lim_{n\to\infty}(U_n - L_n) = 0.
\]

Since $f$ is uniformly continuous on $[0,1]$ by Theorem~\ref{thm:19.2},  
there exists $\delta>0$ such that  
\[
x,y\in[0,1], \ |x-y|<\delta \quad\Rightarrow\quad |f(x)-f(y)|<\varepsilon.
\]
Choosing $N$ such that $\tfrac{1}{N}<\delta$ and $n>N$ gives
\[
0 \le U_n - L_n
  = \frac{1}{n} \sum_{i=0}^{n-1} (M_{i,n} - m_{i,n})
  < \varepsilon.
\]

This proves the claim. The next two theorems show uniformly continuous  
functions have nice properties.
\end{definitionbox}


% ===============================
% 28.1 Definition — Differentiability at a point
% ===============================
\begin{definitionbox}[Definition 28.1 — Differentiability]
\label{def:28.1}
Let $f$ be a real-valued function defined on an open interval that contains a point $a$.  
We say that $f$ is \emph{differentiable at $a$}, or that $f$ \emph{has a derivative at $a$}, if the limit
\[
\lim_{x \to a} \frac{f(x)-f(a)}{x-a}
\]
exists (as a finite real number).
\end{definitionbox}

% ===============================
% 28.2 Theorem — Differentiability implies continuity
% ===============================
\begin{theorembox}[Theorem 28.2]
\label{thm:28.2}
If $f$ is differentiable at a point $a$, then $f$ is continuous at $a$.
\end{theorembox}

% ===============================
% 28.3 Theorem — Algebra rules for derivatives
% ===============================
\begin{theorembox}[Theorem 28.3 — Algebra Rules]
\label{thm:28.3}
Let $f$ and $g$ be functions that are differentiable at a point $a$, and let $c$ be a real constant.  
Then each of the following functions is differentiable at $a$ (whenever it is defined), and their derivatives
are given by:
\begin{enumerate}[(i)]
  \item \textbf{Constant multiple:}
        \[
        (cf)'(a) = c \cdot f'(a).
        \]
  \item \textbf{Sum:}
        \[
        (f+g)'(a) = f'(a) + g'(a).
        \]
  \item \textbf{Product rule:}
        \[
        (fg)'(a) = f(a)g'(a) + f'(a)g(a).
        \]
  \item \textbf{Quotient rule:} If $g(a) \neq 0$, then $f/g$ is differentiable at $a$ and
        \[
        \biggl(\frac{f}{g}\biggr)'(a)
          = \frac{g(a)f'(a) - f(a)g'(a)}{[g(a)]^{2}}.
        \]
\end{enumerate}
\end{theorembox}

% ===============================
% 28.4 Theorem — Chain Rule
% ===============================
\begin{theorembox}[Theorem 28.4 — Chain Rule]
\label{thm:28.4}
Assume $f$ is differentiable at $a$ and $g$ is differentiable at $f(a)$.  
Then the composite function $g \circ f$ is differentiable at $a$, and
\[
(g \circ f)'(a) = g'\bigl(f(a)\bigr)\, f'(a).
\]
\end{theorembox}

\end{document}